% 分野別類題: ガウス積分 解答例
% コンパイル: TEXINPUTS="../../../00_common:$TEXINPUTS" latexmk -xelatex answers.tex
\documentclass[a4paper,11pt]{article}
\usepackage{preamble}
\usepackage{shoakubox}

\begin{document}

\kamokumei{類題演習 解答例 --- ガウス積分}

\daimon{【解法】ガウス積分の解き方と導出}

\noindent
\textbf{(1) 基本公式 $\displaystyle\int_{-\infty}^{\infty} e^{-x^{2}}\,dx = \sqrt{\pi}$ の導出。}
\[
I = \int_{-\infty}^{\infty} e^{-x^{2}}\, dx
\]
とおく。$I > 0$ であり、
\[
I^{2} = \left(\int_{-\infty}^{\infty} e^{-x^{2}}\, dx\right)\left(\int_{-\infty}^{\infty} e^{-y^{2}}\, dy\right)
= \iint_{\mathbb{R}^{2}} e^{-(x^{2}+y^{2})}\, dx\, dy
\]
ここで極座標 $x = r\cos\theta,\ y = r\sin\theta$（$r \ge 0,\ 0 \le \theta < 2\pi$）に変換すると、面積要素は $dx\,dy = r\, dr\, d\theta$ であるから
\[
I^{2} = \int_{0}^{2\pi}\!\! \int_{0}^{\infty} e^{-r^{2}}\, r\, dr\, d\theta
= 2\pi \int_{0}^{\infty} r e^{-r^{2}}\, dr
\]
$u = r^{2}$ とおくと $du = 2r\, dr$ より
\[
\int_{0}^{\infty} r e^{-r^{2}}\, dr = \frac{1}{2}\int_{0}^{\infty} e^{-u}\, du = \frac{1}{2}
\]
したがって $I^{2} = 2\pi \cdot \dfrac{1}{2} = \pi$、$I > 0$ より
\[
\int_{-\infty}^{\infty} e^{-x^{2}}\, dx = \sqrt{\pi}
\]
一般に $a > 0$ に対しては、置換 $t = \sqrt{a}\, x$（$dt = \sqrt{a}\, dx$）により
\[
\int_{-\infty}^{\infty} e^{-a x^{2}}\, dx = \int_{-\infty}^{\infty} e^{-t^{2}}\, \frac{dt}{\sqrt{a}} = \sqrt{\frac{\pi}{a}}
\]

\noindent
\textbf{(2) パラメータ微分による関連積分の導出。}
$\displaystyle F(a) = \int_{-\infty}^{\infty} e^{-a x^{2}}\, dx = \sqrt{\pi}\, a^{-1/2}$ を $a$ の関数とみて微分すると、被積分関数が十分滑らかで積分と微分の順序交換ができるから
\[
F'(a) = \int_{-\infty}^{\infty} (-x^{2}) e^{-a x^{2}}\, dx = -\sqrt{\pi}\cdot\left(-\frac{1}{2}\right) a^{-3/2} = \frac{\sqrt{\pi}}{2} a^{-3/2}
\]
よって
\[
\int_{-\infty}^{\infty} x^{2} e^{-a x^{2}}\, dx = -F'(a) = \frac{\sqrt{\pi}}{2\, a^{3/2}}
\]
のように、高次のモーメントもパラメータ微分を繰り返すことで求められる。

\noindent
\textbf{(3) 平方完成による一般化。}
$a > 0$、$b$ を実数の定数とすると
\[
-a x^{2} + b x = -a\left(x - \frac{b}{2a}\right)^{2} + \frac{b^{2}}{4a}
\]
と平方完成できる。$t = x - \dfrac{b}{2a}$ とおくと（積分域は $x$ 全体を動くとき $t$ も全体を動く）
\[
\int_{-\infty}^{\infty} e^{-a x^{2} + b x}\, dx
= e^{\frac{b^{2}}{4a}} \int_{-\infty}^{\infty} e^{-a t^{2}}\, dt
= \sqrt{\frac{\pi}{a}}\; e^{\frac{b^{2}}{4a}}
\]
となる。

\clearpage
\begin{mondaiboxS}{問1}
次のガウス積分の値を、2重積分と極座標変換を用いて導出せよ。$\displaystyle I = \int_{-\infty}^{\infty} e^{-x^{2}}\, dx$。さらに $a > 0$ に対する $\displaystyle\int_{-\infty}^{\infty} e^{-a x^{2}}\, dx$ の値を求めよ。
\end{mondaiboxS}
\kaitou
【解法】(1) と同じ手順により
\[
I^{2} = \iint_{\mathbb{R}^{2}} e^{-(x^{2}+y^{2})}\, dx\, dy
= \int_{0}^{2\pi}\!\! \int_{0}^{\infty} e^{-r^{2}}\, r\, dr\, d\theta
= 2\pi\left[-\frac{1}{2}e^{-r^{2}}\right]_{0}^{\infty}
= 2\pi \cdot \frac{1}{2} = \pi
\]
$I > 0$ より
\[
\boxed{\; \int_{-\infty}^{\infty} e^{-x^{2}}\, dx = \sqrt{\pi} \;}
\]
続いて $t = \sqrt{a}\, x$ と置換すると $dx = \dfrac{dt}{\sqrt{a}}$ であり、$x \to \pm\infty$ のとき $t \to \pm\infty$ だから
\[
\int_{-\infty}^{\infty} e^{-a x^{2}}\, dx = \frac{1}{\sqrt{a}} \int_{-\infty}^{\infty} e^{-t^{2}}\, dt = \boxed{\; \sqrt{\dfrac{\pi}{a}} \;}
\]
（検算: $a=1$ のとき $\sqrt{\pi/1} = \sqrt{\pi}$ となり最初の結果と一致する。）

\clearpage
\begin{mondaiboxS}{問2}
$a > 0$ とする。$\displaystyle\text{(1)}\ \int_{0}^{\infty} x\, e^{-a x^{2}}\, dx$、$\displaystyle\text{(2)}\ \int_{-\infty}^{\infty} x^{2}\, e^{-a x^{2}}\, dx$ の値を求めよ。
\end{mondaiboxS}
\kaitou
\textbf{(1)} $u = a x^{2}$ とおくと $du = 2ax\, dx$ より
\[
\int_{0}^{\infty} x\, e^{-a x^{2}}\, dx = \frac{1}{2a}\int_{0}^{\infty} e^{-u}\, du = \frac{1}{2a}\Bigl[-e^{-u}\Bigr]_{0}^{\infty} = \boxed{\; \frac{1}{2a} \;}
\]
（検算: $a=1$ のとき値は $1/2$。数値積分と一致することが知られている。）

\smallskip
\textbf{(2)} $\displaystyle F(a) = \int_{-\infty}^{\infty} e^{-a x^{2}}\, dx = \sqrt{\pi}\, a^{-1/2}$ を $a$ で微分すると（積分と微分の順序交換が正当化できる）
\[
F'(a) = -\int_{-\infty}^{\infty} x^{2} e^{-a x^{2}}\, dx
\]
一方
\[
F'(a) = \sqrt{\pi} \cdot \left(-\frac{1}{2}\right) a^{-3/2} = -\frac{\sqrt{\pi}}{2}a^{-3/2}
\]
両辺を比較して
\[
\int_{-\infty}^{\infty} x^{2} e^{-a x^{2}}\, dx = \boxed{\; \frac{\sqrt{\pi}}{2\, a^{3/2}} \;}
\]
（検算: $a=1$ のとき、部分積分 $\displaystyle\int_{-\infty}^{\infty} x^{2}e^{-x^{2}}dx = \left[-\frac{1}{2}x e^{-x^{2}}\right]_{-\infty}^{\infty} + \frac{1}{2}\int_{-\infty}^{\infty} e^{-x^{2}}dx = 0 + \frac{\sqrt{\pi}}{2}$ となり、公式に $a=1$ を代入した値と一致する。）

\clearpage
\begin{mondaiboxS}{問3}
$a > 0$、$b$ を実数の定数とする。平方完成を用いて $\displaystyle\int_{-\infty}^{\infty} e^{-a x^{2} + b x}\, dx$ の値を求めよ。また、この結果を用いて $\displaystyle\int_{-\infty}^{\infty} e^{-x^{2} + 2x}\, dx$ の値を求めよ。
\end{mondaiboxS}
\kaitou
指数部を平方完成すると
\[
-a x^{2} + b x = -a\left(x^{2} - \frac{b}{a}x\right)
= -a\left(x - \frac{b}{2a}\right)^{2} + \frac{b^{2}}{4a}
\]
よって $t = x - \dfrac{b}{2a}$ とおくと（$dt = dx$、積分域は不変）
\[
\int_{-\infty}^{\infty} e^{-a x^{2}+bx}\, dx
= e^{\frac{b^{2}}{4a}} \int_{-\infty}^{\infty} e^{-a t^{2}}\, dt
= \boxed{\; \sqrt{\dfrac{\pi}{a}}\; e^{\frac{b^{2}}{4a}} \;}
\]
$a = 1,\ b = 2$ の場合はこの公式に代入して
\[
\int_{-\infty}^{\infty} e^{-x^{2}+2x}\, dx = \sqrt{\pi}\, e^{\frac{4}{4}} = \boxed{\; \sqrt{\pi}\, e \;}
\]
（検算: $-x^2+2x = -(x-1)^2+1$ であるから、$t=x-1$ とおくと $\int_{-\infty}^{\infty}e^{-(x-1)^2+1}dx = e\int_{-\infty}^{\infty}e^{-t^2}dt = e\sqrt{\pi}$ となり一致する。）

\end{document}
